\begin{definition} \label{definition:rack_of_a_conjugacy_invariant_subset_of_a_group}
Let $G$ be a \CrefAndHyperrefIfExist{definition:group}{group} and $S \subset G$ a \CrefAndHyperrefIfExist{definition:subset_of_a_set}{subset} that is \CrefAndHyperrefIfExist{definition:conjugacy_invariant_subset_of_a_group}{invariant under conjugation}, such that for all $g \in G$ and $s \in S$, $g s g^{-1} \in S$. 
The \hldef{rack of a conjugacy invariant subset} is the \CrefAndHyperrefIfExist{definition:rack}{rack} with underlying set $S$ equipped with the operation $x \rhd y = x y x^{-1}$ for all $x, y \in S$.

In the case that $S = G$, one might call this rack the \hldef{conjugation rack} or \hldef{conjugation quandle} (\Cref{definition:quandle}) of $G$ and denote by \hl{$\operatorname{Conj}(G)$}. Note that the assignment of a \CrefAndHyperrefIfExist{definition:group}{group} to its \CrefAndHyperrefIfExist{definition:conjugation_of_group_elements_and_subgroup_by_a_group_element}{conjugation} rack is \CrefAndHyperrefIfExist{definition:functor_between_categories}{functorial} from the \CrefAndHyperrefIfExist{definition:category_of_groups_of_abelian_groups}{category of groups} to the \CrefAndHyperrefIfExist{definition:category_of_racks}{category of racks}.
\end{definition}